\section{怎样分析词组}

两个或两个以上的词按照一定的方式组合起来，表示比较复杂的意义，但还没有成为句子，这样的一组词叫作词组。比如：“平凡而伟大”、“伟大的平凡”、“非常地伟大”、“平凡得很”，等等，它们既不是意思比较简单的一个词，也不是表达完整意思的一个句子，而是词与词按照一定的组合方式构成的语言单位，即词组。我们要分析词组，就应从分析词与词之间一定的组合方式入手。

怎样分析词与词之间一定的组合方式呢？

1, 分析词与词在构成词组时的先后次序（即词序）。

因为汉语的词序能把词与词的语法关系表示出来；如果词序不同，词组的类别和意义也随之改变。这种情况，在偏正词组、动宾词组、主谓词组中表现得最为明显。例如：“幸福生活”是偏正词组，词与词之间表示修饰和被修饰的关系，表示修饰的词语“幸福”在前（为“偏”），表示被修饰的词语“生活”在后（为“正”），“生活幸福”则是主谓词组，词与词之间表示陈述和被陈述的关系，表示被陈述的词语“生活”在前（叫“主语”），表示陈述的词语“幸福”在后（叫“谓语”）。二者不同的意义和语法关系，通过不同的词序表示出来了，它们之间是不能随意颠倒的。再如：“学哲学”、“反对霸权主义”等，都是动宾词组，词与词之间表示支配与被支配的关系，表示支配关系的词语在前（是“动词”），表示被支配关系的词语在后（是“宾语”）。偏正词组、主谓词组和动宾组词的语法关系纯粹靠词序来表示，这一点必须严格遵守。

即使是词和词平等并列在一起构成联合关系的联合词组，如“工农兵”、“讨论通过”等，各个组成部分次序比较灵活，但也要受到表达内容和语言习惯等条件的限制。比如人们不说“农工兵”或“兵工农”，因为这样既不符合习惯用法，又难于明确地表示出特定的意义。至于“讨论通过”，本是联合词组，如果变动词序，就成了“通过讨论”，在一定的语言环境里，那就由联合词组变成了“介词结构”，意义和语法关系都发生了变化。因此，我们在分析词组时，必须抓住词序，以确定词组的类别和意义。

2. 分析词组，有时还要通过词与词之间的虚词来进行，以求确定词组类别和意义。这些虚词包括连词、介词、副词和助词。词组的组成有时还得靠这些虚词帮助。例如：

“伟大和平凡”、“老师和同学”，通过连词“和”，构成联合关系的联合词组；“伟大的平凡”、“老师的同学”，通过结构助词“的”，构成偏正关系的偏正词组。再如“讨论问题”是表示动宾关系的动宾词组，而“讨论的问题”则通过结构助词“的”，把它变成表示偏正关系的偏正词组，意义当然也就完全不同了。又如：“很生动”、“狠打”是偏正词组，而“生动得很”、“打得狠”则是“述补词组”。这种述补词组的词和词之间有补充说明和被补充说明的关系，前边的是被补充说明的中心词，常由动词或形容词充当；后边的起补充说明作用，叫补语，两部分之间有时就用结构助词“得”来表示语法关系。有人把述补词组归入偏正词组，那么前面的中心词则为“正”，后面的补语则为“偏”了。

一般说来，偏正词组当中有时用上结构助词“的”或“地”；述补词组当中有时用上结构助词“得”；动宾组词的动词和宾语之间有时用上时态助词“了”；联合词组的组成部分之间有时用上连词“和”、“而”、“或”、“并”等，有时用副词“又……又……”联结并列的动词或形容词（如“又高又大”、“又吃又喝”等）。

至于介词，经常附着在名词前面（有时也附着在代词、动词前面），组成“介词结构”（或叫“介词词组”）。介词结构的用途是修饰或者补充说明动词、形容词等，表示时间、处所、方向、对象等。例如：“从上海来”、“为人民服务”中的“从上海”、“为人民”都是介词加名词组成的介词词组——前者修饰“来”，表示处所；后者修饰“服务”，表示对象。另外，助词“的”也可以附在动词或词组后边，合起来成为具有名词功能的“的字结构”。例如：“我的和你的”、“公开的和隐蔽的”、“敢于冲锋陷阵的和心存畏惧临阵脱逃的”等等。“介词结构”和“的字结构”，虽然只是实词和虚词各为一方的组合，而不是实词与实词之间组合而成的词组，但是它们可以同各种词组套用在一起，成为比较复杂的词组。

3. 对于比较复杂的词组，可以采用不断切分的方法，分析其中词语之间的层次和关系。复杂的词组是一层套一层的，可以不断切分。这种逐层的切分，可以用图解法来表示。例如：

\begin{enumerate}[label=\textcircled{\arabic*}]
\item 革命的友谊和战斗的团结

\begin{dependency}[theme = simple, edge unit distance=1em]
	\begin{deptext}[column sep=1em]
		革命 \& 的 \& 友谊 \& 和 \& 战斗 \& 的 \& 团结 \\
	\end{deptext}
	\depedge{1}{3}{偏正}
	\depedge{5}{7}{偏正}
	\depedge{1}{7}{联合}
\end{dependency}
\newpage

\item 语言生动的文章

\begin{dependency}[theme = simple, edge unit distance=0.5em]
	\begin{deptext}[column sep=0.5em]
		语言 \& 生动 \& 的 \& 文章 \\
	\end{deptext}
	\depedge{1}{2}{主谓}
	\depedge{1}{4}{偏正}
\end{dependency}

\item 提高分析问题和解决问题的能力

\begin{dependency}[theme = simple, edge unit distance=0.5em]
	\begin{deptext}[column sep=0.5em]
		提高 \& 分析 \& 问题 \& 和 \& 解决 \& 问题 \& 的 \&能力 \\
	\end{deptext}
	\depedge{2}{3}{动宾}
	\depedge{5}{6}{动宾}
	\depedge{2}{6}{联合}
	\depedge{2}{8}{偏正}
	\depedge{1}{8}{动宾}
\end{dependency}

\item 写得生动活泼又细腻深刻

\begin{dependency}[theme = simple, edge unit distance=0.5em]
	\begin{deptext}[column sep=0.5em]
		写 \& 得 \& 生动 \& 活泼 \& 又 \& 细腻 \& 深刻\\
	\end{deptext}
	\depedge{3}{4}{联合}
	\depedge{6}{7}{联合}
	\depedge{3}{7}{联合}
	\depedge{1}{7}{述补}
\end{dependency}

\item 对于顽抗到底的统统消灭

\begin{dependency}[theme = simple, edge unit distance=0.5em]
	\begin{deptext}[column sep=0.5em]
		对 \& 于 \& 顽抗 \& 到底 \& 的 \& 统统 \& 消灭\\
	\end{deptext}
	\depedge{3}{4}{述补}
	\depedge{3}{6}{的字结构}
	\depedge{1}{6}{介词结构}
	\depedge{1}{7}{偏正}
\end{dependency}
\newpage

\item 为实现四个现代化的宏伟目标而奋斗

\begin{dependency}[theme = simple, edge unit distance=0.5em]
	\begin{deptext}[column sep=0.5em]
		为 \& 实现 \& 四个 \& 现代化 \&  的 \& 宏伟 \& 目标 \& 而 \& 奋斗\\
	\end{deptext}
	\depedge{3}{4}{偏正}
	\depedge{6}{7}{偏正}
	\depedge{3}{7}{偏正}
	\depedge{2}{7}{动宾}
	\depedge{1}{7}{介词结构}
	\depedge{1}{9}{偏正}
\end{dependency}
\end{enumerate}

4. 运用分析词组的常识和方法，可以把句子写得通畅整齐，也可以避免和改正由于词语和词语之间的次序安排不当而造成的病句。例如：

\begin{enumerate}[label=\textcircled{\arabic*}]
\item 这些邮包不能放颠倒、互相碰。

（“放颠倒”是述补词组，“互相碰”是偏正词组，读起来不整齐，不顺口。可把“放颠倒”改为“颠倒放”，就整齐顺口了。）

\item 这篇文章有充实的内容，立意新颖。

（把动宾词组“有充实的内容”改为主谓词组“内容充实”，与后面的主谓词组“立意新颖”相配，既整齐顺口，又较简练。）

\item “干”字是一个很简单、普通得很的字。

（把述补词组“普通得很”改为偏正词组“很普通”，与前面的偏正词组“很简单”相配，就整齐顺口了。）
\end{enumerate}
\newpage

再如：

\begin{enumerate}[label=\textcircled{\arabic*}]
\item 我要做一切党希望做的事情。

（“一切党”是个偏正词组，表示“所有的党”，它与“希望做”构成一个比较复杂的主谓词组，来修饰“事情”。而这却与作者要表达的意思大有出入。应该将“一切”调到“事情”之前，与“事情”组成偏正词组，意思就既明确又正确了。）

\item 这部作品揭示了人民战争取得真正胜利的原因。

（“真正胜利”是个偏正词组，充当“取得”的宾语，与“取得”组成比较复杂的动宾词组；它们一起与“人民战争”组成比较复杂的主谓词组，去修饰“原因”。然而，作者的本意并非在强调与“虚假胜利”相对立的“真正胜利”，而是在强调作品揭示了与“虚假原因”相对立的“真正原因”。应该将“真正”调到“原因”之前，与“原因”组成偏正词组，意思就既明确又正确了。）

\item 在社会主义建设事业中，他们发挥了无穷蕴藏着的力量。

（“无穷蕴藏”是个偏正词组，强调的是“蕴藏”“没有穷尽、没有限度”；而作者的本意是强调“力量”“没有穷尽、没有限度”。应该将“无穷”调到“力量”之前，与“力量”组成偏正词组，充当“发挥”的宾语，意思就既明确又正确了。）
\end{enumerate}

至于成语和专有名称（一种特定对象的名称，指人名、地名、机关团体名之类），由于组成成分固定，结构紧密，其成分、结构不能随意增减变换。例如成语“雪中送炭”不能写成“送炭雪中”或“霜中送炭”，“车水马龙”不能写成“车龙马水”或“车马水龙”，等等，我们就叫它固定词组。再如专有名称“郭沫若”、“乌鲁木齐”、“狂人日记”、“北京图书馆”、“中华人民共和国”，等等，都不能随意改动。这些成语或专有名称在句子中的作用常常相当于一个词，充当句子的一个成分。在分析句子时，一般就不再进 行内部结构的分析了。
\newpage